\section{Introduction}

The supply of machine intelligence is expanding rapidly. Models differ not
only in overall quality, but in scientific knowledge, formal reasoning, code
generation, critique, tool use, latency, and price. Yet most inference systems
still consume this supply through one of two simple patterns. A request is sent
to one general-purpose frontier model, or it is broadcast to several strong
models whose answers are subsequently ranked, voted on, or synthesized.

Both patterns place the model at the center of the system. They ask which model
should answer, or how many answers should be generated. They do not directly
ask the more consequential question: \emph{what is unresolved in the current
task state, and what computation is most likely to resolve it?}

This distinction matters because model calls are not independent units of
intelligence. Two highly capable models may reproduce the same reasoning and
the same error. A smaller specialist may expose a decisive counterexample that
a second frontier model misses. In a long-running workflow, no language model
can determine whether a test passed without observing the environment. And
after enough evidence has accumulated, an additional answer can add cost and
latency without changing the decision.

We introduce \systemname, a model system built around \emph{allocation
intelligence}: the ability to allocate heterogeneous model capability according
to the marginal value of evidence in the current task state. \systemname does
not treat a model output as a vote or an action. It treats each model--role call
as a conditional evidence source and preserves one authoritative path from the
canonical task state to the next response or external action.

This perspective produces a system that is asymmetric by design. Cheap and
specialized models can explore, solve, or critique; programmatic checks and
tools can contribute scoped observations; and a stronger model can retain
decision authority without performing every unit of work itself. Compute
efficiency is therefore not a secondary commercial constraint. It is evidence
that the system can distinguish useful computation from redundant computation.

\paragraph{A new scaling dimension.}
Model progress increases the intelligence available to the system. Allocation
progress increases how effectively that intelligence is converted into task
outcomes. \systemname benefits from both. Our aim is to make system capability
compound faster than frontier-compute consumption, not by weakening the final
decision-maker, but by spending expensive intelligence only where it has high
marginal value.

\paragraph{Contributions.}
This paper makes four claims.
\begin{enumerate}
  \item We formulate multi-model inference as adaptive evidence acquisition
  over a changing task state, rather than as static model selection or answer
  aggregation.
  \item We present a system principle---\emph{one state, many perspectives,
  one commit}---that separates plural evidence generation from singular answer
  and action authority.
  \item We describe an uncertainty-directed architecture that jointly reasons
  about task structure, workflow phase, model capability, evidence readiness,
  cost, and risk.
  \item We define an empirical and learning agenda for allocation intelligence:
  performance--compute frontiers, conditional contribution, long-horizon task
  success, and outcome-driven improvement of the allocation policy.
\end{enumerate}

